\section{The process and the subtraction boundary}\label{sec:model}

Let $S$ be a finite subset of $\N_0$ containing the root $0$.  If
$S\ne\{0\}$, select $v\in S\setminus\{0\}$ uniformly and set
\begin{equation}
 C_v(S)=(S\setminus\{v\})\cup\{v-1\}.                 \label{eq:update}
\end{equation}
If $v-1$ is already occupied, the two piles coalesce; the state records no
multiplicity.  Let $T_S$ be the number of updates until $\{0\}$ is reached.
This is an \emph{active-pile} chain: uniform selection is among the current
nonroot piles.  Uniform selection from a fixed geometric set with lazy empty
events gives a different discrete-time law.

Coalescing random walks and their voter-model duals form a classical
interacting-particle framework; full occupancy, consensus time, and
meeting-time comparison are established themes
\cite{Cox1989,Cooper2013,Kanade2023}.  Graphical constructions and recurrence
questions on infinite or rooted carriers are likewise standard
\cite{Benjamini2016}.  Exact one-dimensional coalescence and reductions to
simple-walk first passage also have a substantial history
\cite{Ermakov1997}.  We assign zero contribution credit to coalescing-walk
formalism, graphical arrows, voter duality, meeting/hitting comparisons,
ballot identities, and central-binomial asymptotics.

Three closer interfaces sharpen that subtraction.  \citet{Assiotis2018}
develops stochastic coalescing flows associated with birth--death chains,
while \citet{SniadyUrban2026} treat prescribed coalescence patterns for
nearest-neighbor systems, including birth--death chains, using ordered
interval labels and an explicit finite coalescing construction.
\citet[Theorem~3]{HitczenkoWesolowski2025} identify, for step-initial
TASEP, the expected active-particle count as the derivative of the expected
total jump count.  Their kernels and observables differ from ours, but the
site-arrow flow, ordered-label-block, coalescence-pattern, and generic
active-count--jump-current mechanisms all receive zero contribution credit
here.

The literal chain in \eqref{eq:update} differs in three linked respects.
Motion is deterministic and one-way, the scheduler is uniform among current
piles, and $T_S$ counts embedded updates rather than elapsed coalescence time.
The residual question is whether this update count admits an exact formula
from an arbitrary initial set.  A bounded primary-source search found no
statement of the theorem package below.  That non-hit is not a novelty
certificate, so external release remains on hold.

The internal subtraction is also strict.  P114 already owns synchronous
rooted-forest peeling.  P117 acts on labelled cyclic binary words by flipping
every odd maximal run; its boundary-parity eroder has neither piles, a root,
nor a random scheduler.  P121 instead selects a current adjacent separator,
performs a product-plus-one merge, and obtains a random-BST/Yule deletion
history.  Here the selected object is a current pile, a move need not
coalesce, and the observable is the number of embedded effective updates.
P126 concerns synchronous length-increasing refinement.  Rooted
monotonicity, eroder language, generic coalescence, adjacent interfaces, and
random scheduling therefore receive no internal credit.  The only claimed
residual is the exact arbitrary-state embedded-update mean for the literal
deterministic-rootward, uniform-active-pile chain, together with its support
and full-start consequences.
